\section{Three connected research questions}
In this era, mobile phones have taken a large amount of time from us. Applications on the mobile phones are competing for users’ attention and time. To achieve this purpose, applications sometimes adopt addictive recommendation algorithms, which is harmful for users’ mental health.

Jean Tirole was awarded the Nobel Prize in 2014 for his analysis on how regulators make rules to make large companies benefit their users \cite{nobel2014}. When it comes to applications’ recommendation algorithms, regulators can also influence applications, driving applications to adopt recommendation algorithms that are healthier for the user.

\textbf{Q1—Economics:} How do external rewards from regulators shape applications’ recommendation algorithms?\par
\textbf{Q2—Computation:} How to evaluate whether a recommendation is healthy or harmful?\par
\textbf{Q3—Behavior:} Under what circumstances will the user choose to use applications that adopt healthier recommendation algorithms?\par

The three questions focus on applications’ recommendation algorithms. The goal is to make applications provide healthier algorithms for users. To achieve this goal, as the teaser figure shows, we need to firstly define what a healthy algorithm is (Question 2). Many studies focus on how to make recommendation algorithms that keep users using \cite{zou2019}. Conversely, we want to make applications willing to adopt healthy recommendation algorithms by designing an external reward (Question 1). Allcott et al.~\cite{allcott2022} claim that people do not intentionally spend so much time on mobile phones \cite{wu2026}. This indicates mobile phone usage behaviors sometimes diverge from users’ rational intentions. The goal of making applications provide healthier algorithms for users cannot be achieved unless users intentionally choose to use those applications that adopt healthier recommendation algorithms (Question 3).

\section{Answer to Q1:}
My hypothesis is applications will adopt healthier recommendation algorithms if the reward for healthy algorithms is big enough. Players in this model are two applications that aim to attract users’ attention and time. Each player chooses to cooperate (C) by adopting healthy algorithms or chooses to defect (D) by adopting addictive algorithms. The regulator will award the application who chooses to cooperate a reward r. The hypothetical payoff is (3+r, 3+r) for CC, (4, 1+r) for DC, (1+r, 4) for CD, and (2, 2) for DD. In this case, (a, b) means application A’s payoff is a and application B’s payoff is b; CD means application A chooses to cooperate and application B chooses to defect.

My solution concept is Nash equilibrium \cite{nash1950}. When the reward from the regulator is big enough, both applications will tend to choose to cooperate by adopting healthier recommendation algorithms, no matter what their competitor’s choice is. In this case, when r is larger than 1, both applications will tend to choose cooperation no matter what their competitor’s choice is, since choosing cooperation brings a bigger payoff.

\section{Answer to Q2:}
In this proposal, “healthy recommendation algorithm” is a recommendation algorithm that makes users feel that they have benefited from the recommended content instead of feeling regretful for interacting with the recommended content. As the teaser figure shows, to evaluate whether a recommendation algorithm is healthy or unhealthy, users are asked to rate recommended content. Based on users’ feedback on whether they benefit from the content or regret watching the content, the applications’ recommendation algorithms are labeled as “healthy” or “unhealthy”. As a preliminary design: If the number of users who benefit from the recommended content is more than the number of users who regret interacting with the recommended content, the corresponding recommendation algorithm is labeled as “healthy”. Otherwise, the corresponding recommendation algorithm is labeled as “unhealthy”.


\balance
\section{Answer to Q3:}
One explanation is the user will choose to use the application that adopts a healthy recommendation algorithm as long as the user knows which application adopts the healthy algorithm. The competing explanation is the user will choose to use the application that provides content that attracts the user, no matter whether its recommendation algorithm is healthy. In the proposed experiment, two applications are shown for the users to choose from. One application adopts a healthy recommendation algorithm, while the other adopts an unhealthy recommendation algorithm. Users are told which application adopts a healthy algorithm before they select an application. After the users make their choice, they will be asked about the reason for their choice. If users report that their choice is based on whether the application adopts a healthy recommendation algorithm, then it provides preliminary evidence for the explanation “the user will choose to use the application that adopts a healthy recommendation algorithm as long as the user knows which application adopts the healthy algorithm”. However, self-reported results alone are not solid enough to claim which explanation holds.

\section{Advanced development and future directions}

\begin{table}[h]\caption{From laureate foundations to a new interdisciplinary contribution.}
\begin{tabularx}{\columnwidth}{@{}p{1.6cm}Y@{}}\toprule
\textbf{Horizon} & \textbf{Milestone and evidence}\\
\midrule
Foundation & Jean Tirole was awarded the Nobel Prize in 2014 for his analysis on how regulators make rules to make large companies benefit their users \cite{nobel2014}. This indicates that how to make competition benefit the public has been a long-standing and important question.\\
PS1 & Applications’ competition for users’ attention and time is becoming increasingly intense.\\
Next study & My proposal advances research on this question. This proposal integrates economics, computer science and behavioral science. It focuses on how regulators’ reward shapes applications’ recommendation algorithms. It also studies human behavior when users face different recommendation algorithms.\\
Longer term & Hence, it provides a more comprehensive perspective on how to make applications provide healthier algorithms for users.\\
2056 & In 2056, I hope to be awarded the Nobel Prize for combining computer science, economics and behavioral science to address the problem of how to make applications adopt healthier recommendation algorithms.\\
\bottomrule\end{tabularx}\end{table}

\noindent\textbf{Expected 2056 Nobel Prize and Turing Award contribution statement.}
By 2056, I aspire to be recognized for the method addressing the question of how to make applications provide healthy algorithms in the era when applications compete for users’ attention and time. Building on Jean Tirole’s analysis on how regulators make rules to make large companies benefit their users \cite{nobel2014}, I will integrate the method to make applications adopt healthy recommendation algorithms with external reward, the method to distinguish healthy and unhealthy recommendation algorithms and users’ choice facing different recommendation algorithms to advance research on the problem of how to make competitors cooperate and benefit society. Progress will be demonstrated by results of experiments shown in Answer to Q1, Q2 and Q3, with intellectual merits in economics, computer science and behavioral science and practical impacts for netizens.

\subsection*{Open Science Statement}
GitHub Code repository: {\def\UrlBreaks{\do a\do b\do c\do d\do e\do f\do g\do h\do i\do j\do k\do l\do m\do n\do o\do p\do q\do r\do s\do t\do u\do v\do w\do x\do y\do z\do A\do B\do C\do D\do E\do F\do G\do H\do I\do J\do K\do L\do M\do N\do O\do P\do Q\do R\do S\do T\do U\do V\do W\do X\do Y\do Z\do 0\do 1\do 2\do 3\do 4\do 5\do 6\do 7\do 8\do 9\do /\do -\do _\do .\do ?\do =\do &}\url{https://github.com/yw650/Recommendation-Algorithms}}

Google Colab: {\def\UrlBreaks{\do a\do b\do c\do d\do e\do f\do g\do h\do i\do j\do k\do l\do m\do n\do o\do p\do q\do r\do s\do t\do u\do v\do w\do x\do y\do z\do A\do B\do C\do D\do E\do F\do G\do H\do I\do J\do K\do L\do M\do N\do O\do P\do Q\do R\do S\do T\do U\do V\do W\do X\do Y\do Z\do 0\do 1\do 2\do 3\do 4\do 5\do 6\do 7\do 8\do 9\do /\do -\do _\do .\do ?\do =\do &}\url{https://colab.research.google.com/drive/1mYoomco3ywsZb8PaX2oh-wjGZkXDkRil?usp=sharing}}

Submitted code commit: {\def\UrlBreaks{\do a\do b\do c\do d\do e\do f\do g\do h\do i\do j\do k\do l\do m\do n\do o\do p\do q\do r\do s\do t\do u\do v\do w\do x\do y\do z\do A\do B\do C\do D\do E\do F\do G\do H\do I\do J\do K\do L\do M\do N\do O\do P\do Q\do R\do S\do T\do U\do V\do W\do X\do Y\do Z\do 0\do 1\do 2\do 3\do 4\do 5\do 6\do 7\do 8\do 9\do /\do -\do _\do .\do ?\do =\do &}\nolinkurl{7720aa975f87b6512bfc5d2852f3b4ee5179e1c3}}

Paper source: {\def\UrlBreaks{\do a\do b\do c\do d\do e\do f\do g\do h\do i\do j\do k\do l\do m\do n\do o\do p\do q\do r\do s\do t\do u\do v\do w\do x\do y\do z\do A\do B\do C\do D\do E\do F\do G\do H\do I\do J\do K\do L\do M\do N\do O\do P\do Q\do R\do S\do T\do U\do V\do W\do X\do Y\do Z\do 0\do 1\do 2\do 3\do 4\do 5\do 6\do 7\do 8\do 9\do /\do -\do _\do .\do ?\do =\do &}\url{https://github.com/yw650/PS1-Yongzhuo}}

An IPYNB file is included in the GitHub repository. Open and run it.

\subsection*{Statement of Contribution to the UN’s SDGs}
Hugging Face: {\def\UrlBreaks{\do a\do b\do c\do d\do e\do f\do g\do h\do i\do j\do k\do l\do m\do n\do o\do p\do q\do r\do s\do t\do u\do v\do w\do x\do y\do z\do A\do B\do C\do D\do E\do F\do G\do H\do I\do J\do K\do L\do M\do N\do O\do P\do Q\do R\do S\do T\do U\do V\do W\do X\do Y\do Z\do 0\do 1\do 2\do 3\do 4\do 5\do 6\do 7\do 8\do 9\do /\do -\do _\do .\do ?\do =\do &}\url{https://huggingface.co/spaces/dku-comsci-econ206-2026/Recommendation_Algorithms}}

Through an open-access online interactive game, this project aims to support SDG 4—Quality Education by helping learners learn how external reward from regulators would shape different applications’ recommendation algorithms through interacting with the online game. Interact with this online game and monitor the changes.

\subsection*{Submission milestones}
First draft: September 13, 2026, 11:00 P.M. Beijing time (UTC+8). Class review: September 14. Proposed: rebuttal September 16 and revised submission September 20, both at 11:00 P.M.
